% =============================================================================
%                  WEEK 3: REPEATED AND EXTENSIVE-FORM GAMES
% =============================================================================
\section{Week 3: Repeated and Extensive-Form Games}

% -----------------------------------------------------------------------------
%                              3.1 TOPIC SUMMARY
% -----------------------------------------------------------------------------
\subsection{Topic Summary}

\begin{keybox}[Learning Objectives]
By the end of this week, you will be able to:
\begin{enumerate}
  \item Define the core concepts of \textbf{extensive form games}, including game trees, choice nodes, and terminal nodes.
  \item Formulate real-world scenarios as \textbf{extensive-form games} with perfect or imperfect information.
  \item Identify \textbf{pure strategies} in extensive form games as complete plans specifying actions at every choice node.
  \item Derive the \textbf{induced normal form} from an extensive form game and analyse its Nash equilibria.
  \item Define and compute \textbf{subgame-perfect equilibria (SPE)} using backward induction.
  \item Distinguish between \textbf{finitely} and \textbf{infinitely repeated games} and their equilibrium properties.
  \item Analyse cooperation-sustaining strategies (\textbf{Tit-for-Tat}, \textbf{Trigger}) in infinitely repeated Prisoner's Dilemma.
\end{enumerate}
\end{keybox}

\separator

\subsubsection*{The Big Picture}

This week moves beyond normal form (simultaneous) games to settings where players act \textbf{sequentially} or \textbf{repeatedly}. Two key representations are introduced:

\begin{itemize}
  \item \textbf{Extensive form games:} Game trees that model sequential moves, with \textbf{perfect information} (all history visible) or \textbf{imperfect information} (some moves hidden via information sets). The key refinement is \textbf{subgame-perfect equilibrium}, which eliminates non-credible threats.
  \item \textbf{Repeated games:} The same stage game played multiple times. Finitely repeated games collapse to stage-game NE by backward induction; \textbf{infinitely repeated games} can sustain cooperation via strategies like Tit-for-Tat and Trigger, provided the discount factor is sufficiently large.
\end{itemize}

\separator

\subsubsection*{What Examiners Like to Ask}

\begin{itemize}
  \item Write the \textbf{formal definition} of a perfect-information extensive form game (the 8-tuple).
  \item Given a game tree, enumerate all \textbf{pure strategies} for each player.
  \item Derive the \textbf{induced normal form} (payoff matrix) from a game tree.
  \item Find all \textbf{Nash equilibria} of the induced normal form, then identify which are \textbf{subgame-perfect}.
  \item Apply \textbf{backward induction} to compute SPE.
  \item Explain why every perfect-information game has a \textbf{pure strategy NE}.
  \item Show that the only NE of a \textbf{finitely repeated} Prisoner's Dilemma is always Defect.
  \item For infinitely repeated PD: compute the \textbf{minimum discount factor} $\beta$ for Tit-for-Tat or Trigger to be a NE.
  \item Define \textbf{information sets} and explain the difference between perfect and imperfect information games.
\end{itemize}

\separator

% -----------------------------------------------------------------------------
%           3.2 OVERVIEW + CORE CONCEPTS + EXTENDED THEORY
% -----------------------------------------------------------------------------
\subsection{Overview + Core Concepts + Extended Theory}

\subsubsection*{Roadmap}
\begin{enumerate}
  \item Extensive form games: motivation and applications
  \item Perfect-information games: formal definition
  \item Strategies in extensive form games
  \item Induced normal form and its redundancy
  \item Pure strategy NE existence theorem
  \item Subgame-perfect equilibrium and backward induction
  \item Imperfect-information games and information sets
  \item Repeated games: finite and infinite
  \item Payoff definitions: average reward and discounted reward
  \item Cooperation strategies: Tit-for-Tat and Trigger
\end{enumerate}

\bigseparator

% --- Motivation ---
\subsubsection{From Normal Form to Extensive Form}

\begin{defbox}[Why Extensive Form?]
The normal form (payoff matrix) is universal but has limitations:
\begin{itemize}
  \item Assumes \textbf{simultaneous} moves---cannot capture sequential structure
  \item Can be \textbf{exponentially larger} than the equivalent extensive form
  \item Misses solution concepts that depend on the \textbf{order of play} (e.g., subgame perfection)
\end{itemize}

Extensive form games are visualised as \textbf{game trees} and are widely used in AI:
\begin{center}
\renewcommand{\arraystretch}{1.3}
\begin{tabular}{@{} l p{8.5cm} @{}}
\toprule
\textbf{Application} & \textbf{How Extensive Form Is Used} \\
\midrule
Chess, Go & Sequential moves modelled explicitly; minimax / alpha-beta pruning \\
Poker & Imperfect information; extensive-form fictitious self-play \\
Automated negotiation & Multi-stage strategic interaction with bluffing \\
Sequential auctions & Bids placed in turn; captures information and tactics \\
Security / patrolling & Sequential actions to maximise system utility (e.g., Stackelberg games) \\
Multi-agent RL & Sequential, stage-dependent strategies with private histories \\
\bottomrule
\end{tabular}
\end{center}
\end{defbox}

\separator

% --- Perfect-Information Games ---
\subsubsection{Perfect-Information Extensive Form Games}

\begin{defbox}[Perfect-Information Game (Formal Definition)]
A \textbf{finite perfect-information game} in extensive form is a tuple:
\[
G = (N,\; A,\; H,\; Z,\; \chi,\; \rho,\; \sigma,\; u)
\]
where:
\begin{itemize}
  \item $N$: set of $n$ players
  \item $A$: set of actions
  \item $H$: set of \textbf{nonterminal choice nodes}
  \item $Z$: set of \textbf{terminal nodes} (disjoint from $H$)
  \item $\chi: H \to 2^A$: \textbf{action function}---assigns available actions to each choice node
  \item $\rho: H \to N$: \textbf{player function}---assigns a player to each choice node
  \item $\sigma: H \times A \to H \cup Z$: \textbf{successor function}---maps (node, action) to the next node; injective (if $\sigma(h_1, a_1) = \sigma(h_2, a_2)$ then $h_1 = h_2$ and $a_1 = a_2$)
  \item $u = (u_1, \ldots, u_n)$: \textbf{utility functions} $u_i: Z \to \mathbb{R}$
\end{itemize}
\end{defbox}

\begin{tipbox}[Exam Tip: Tree Terminology]
Since choice nodes form a tree:
\begin{itemize}
  \item Each node can be identified with its \textbf{history}---the sequence of actions from the root
  \item The \textbf{descendants} of node $h$ are all nodes in the subtree rooted at $h$
\end{itemize}
\end{tipbox}

\separator

% --- Strategies ---
\subsubsection{Strategies in Perfect-Information Games}

\begin{defbox}[Pure Strategy (Extensive Form)]
A \textbf{pure strategy} for player $i$ is a \textit{complete} specification of which action to take at \textbf{every} choice node belonging to $i$:
\[
s_i \in \prod_{h \in H,\; \rho(h) = i} \chi(h)
\]

\textbf{Key point:} A strategy must specify an action at every node of player $i$, \textit{even nodes that cannot be reached} given the other choices in the strategy.
\end{defbox}

\begin{exbox}[Example: Sharing Game]
\textbf{Setup:} A brother and sister share two identical presents. The brother proposes a split ($2$--$0$, $1$--$1$, or $0$--$2$), then the sister accepts or rejects. Rejection gives both $0$.

\textbf{Strategies:}
\begin{itemize}
  \item Player 1 (brother): 3 strategies: $\{2\text{--}0,\; 1\text{--}1,\; 0\text{--}2\}$
  \item Player 2 (sister): $2^3 = 8$ strategies (accept/reject at each of 3 nodes):

  $S_2 = \{(\text{Y,Y,Y}), (\text{Y,Y,N}), (\text{Y,N,Y}), \ldots, (\text{N,N,N})\}$
\end{itemize}
\end{exbox}

\begin{exbox}[Example: Two-Player Game Tree with Induced Normal Form]
\textbf{Players:} 1 has choice nodes with actions $\{A,B\}$ and $\{G,H\}$; 2 has choice nodes with actions $\{C,D\}$ and $\{E,F\}$.

\textbf{Strategy sets:}
$S_1 = \{(A,G),\; (A,H),\; (B,G),\; (B,H)\}$, \quad $S_2 = \{(C,E),\; (C,F),\; (D,E),\; (D,F)\}$

\textbf{Induced normal form:}
\begin{center}
\renewcommand{\arraystretch}{1.3}
\begin{tabular}{@{} c | c c c c @{}}
 & $(C,E)$ & $(C,F)$ & $(D,E)$ & $(D,F)$ \\
\midrule
$(A,G)$ & $3,8$ & $3,8$ & $8,3$ & $8,3$ \\
$(A,H)$ & $3,8$ & $3,8$ & $8,3$ & $8,3$ \\
$(B,G)$ & $5,5$ & $2,10$ & $5,5$ & $2,10$ \\
$(B,H)$ & $5,5$ & $1,0$ & $5,5$ & $1,0$ \\
\end{tabular}
\end{center}

\textbf{Redundancy:} Only 5 distinct outcomes in the tree, but 16 cells in the matrix---many duplicated because unreachable nodes don't affect payoffs.
\end{exbox}

\separator

% --- Extensive vs Normal Form ---
\subsubsection{Extensive Form vs.\ Normal Form}

\begin{thmbox}[Relationship Between Representations]
\begin{itemize}
  \item \textbf{Every} perfect-information extensive form game has a corresponding normal form game (the induced normal form).
  \item The reverse is \textbf{not} always possible---not every normal form game can be represented as a perfect-information extensive form game.
  \item \textbf{Example:} Prisoner's Dilemma cannot be represented as a perfect-information game (simultaneous moves require imperfect information).
\end{itemize}
\end{thmbox}

\begin{thmbox}[Pure Strategy NE Existence]
Every finite perfect-information game in extensive form has a \textbf{pure strategy Nash equilibrium}.

\textbf{Proof idea:} Since players take turns and observe all prior moves, randomisation is never needed---backward induction yields a pure-strategy equilibrium.

\textbf{Important:} This does \textbf{not} hold for imperfect-information games (e.g., Matching Pennies has no pure NE).
\end{thmbox}

\bigseparator

% --- Subgame-Perfect Equilibrium ---
\subsubsection{Subgame-Perfect Equilibrium}

\begin{defbox}[Subgame]
Given a perfect-information game $G$, the \textbf{subgame} of $G$ rooted at node $h$ is the restriction of $G$ to the descendants of $h$.

The set of subgames of $G$ consists of all subgames rooted at any node in $G$ (including $G$ itself).
\end{defbox}

\begin{defbox}[Subgame-Perfect Equilibrium (SPE)]
A strategy profile $s$ is a \textbf{subgame-perfect equilibrium} of $G$ if, for every subgame $G'$ of $G$, the restriction of $s$ to $G'$ is a Nash equilibrium of $G'$.

\textbf{Key properties:}
\begin{itemize}
  \item Every SPE is a Nash equilibrium (but not vice versa)
  \item SPE eliminates equilibria based on \textbf{non-credible threats}
  \item Every perfect-information game has at least one SPE
\end{itemize}
\end{defbox}

\begin{tipbox}[Exam Tip: SPE vs NE]
If a question asks ``find all Nash equilibria, then identify which are subgame-perfect'':
\begin{enumerate}
  \item Find all NE from the induced normal form
  \item Check each NE against every subgame
  \item Discard any NE where a player's strategy is not a best response in some subgame (non-credible threat)
\end{enumerate}
\end{tipbox}

\begin{exbox}[Example: Filtering NE to SPE]
From the game tree example above, the induced normal form has three Nash equilibria:
\begin{itemize}
  \item $\{(A,G),\; (C,F)\}$, \quad $\{(A,H),\; (C,F)\}$, \quad $\{(B,H),\; (C,E)\}$
\end{itemize}

\textbf{Only $\{(B,H),\; (C,E)\}$ is subgame-perfect}---the other two involve non-credible threats (actions at unreached nodes that are not best responses in their respective subgames).
\end{exbox}

\separator

% --- Backward Induction ---
\subsubsection{Computing SPE: Backward Induction}

\begin{thmbox}[Backward Induction Algorithm]
\textbf{Procedure:}
\begin{enumerate}
  \item Start at the \textbf{bottom-most} (deepest) subgames of the tree
  \item At each terminal decision node, the player chooses the action maximising their payoff
  \item Replace the subtree with the resulting payoff vector
  \item Move up one level and repeat
\end{enumerate}

\textbf{Properties:}
\begin{itemize}
  \item Guarantees finding a \textbf{subgame-perfect} equilibrium
  \item Implemented as a single \textbf{depth-first traversal}---runs in time \textbf{linear} in the size of the game tree
  \item Avoids non-credible threats by construction
\end{itemize}
\end{thmbox}

\begin{exbox}[Centipede Game]
In the Centipede game, backward induction predicts that the first player \textbf{immediately takes} (defects) at the first node, even though continued cooperation would yield higher payoffs. This is a classic case where the SPE prediction seems counterintuitive---in experiments, players often cooperate for several rounds.
\end{exbox}

\bigseparator

% --- Imperfect Information ---
\subsubsection{Imperfect-Information Games}

\begin{defbox}[Imperfect-Information Game]
A \textbf{finite imperfect-information game} extends the perfect-information tuple with \textbf{information sets}:
\[
G = (N,\; A,\; H,\; Z,\; \chi,\; \rho,\; \sigma,\; u,\; I)
\]
where:
\begin{itemize}
  \item $(N, A, H, Z, \chi, \rho, \sigma, u)$ is a perfect-information game
  \item $I = (I_1, \ldots, I_n)$, where $I_i = (I_{i,1}, \ldots, I_{i,k_i})$ is a partition of player $i$'s choice nodes such that for any two nodes $h, h'$ in the same information set $I_{i,j}$:
  \begin{itemize}
    \item $\chi(h) = \chi(h')$ \quad (same available actions)
    \item $\rho(h) = \rho(h')$ \quad (same player)
  \end{itemize}
\end{itemize}

\textbf{Interpretation:} If $h$ and $h'$ are in the same information set, the player \textbf{cannot distinguish} between them---they don't know which node they're at.
\end{defbox}

\begin{defbox}[Pure Strategy (Imperfect Information)]
A pure strategy for player $i$ selects one action in each \textbf{information set} (not each node):
\[
s_i \in \prod_{I_{i,j} \in I_i} \chi(I_{i,j})
\]

Perfect-information games are the special case where every information set is a \textbf{singleton}.
\end{defbox}

\begin{tipbox}[Key Relationships]
\begin{center}
\renewcommand{\arraystretch}{1.3}
\begin{tabular}{@{} l p{9cm} @{}}
\toprule
\textbf{Direction} & \textbf{Result} \\
\midrule
Perfect info $\to$ normal form & Always possible (induced normal form) \\
Normal form $\to$ perfect info & \textbf{Not} always possible (e.g., Prisoner's Dilemma) \\
Normal form $\to$ imperfect info & Always possible (trivially: one information set per player) \\
\bottomrule
\end{tabular}
\end{center}

\textbf{Consequence:} The set of perfect-information games is \textbf{strictly smaller} than the set of all normal form games. Imperfect-information games are equally expressive as normal form games.

Imperfect-information games may \textbf{not} always have a pure strategy NE (unlike perfect-information games).
\end{tipbox}

\bigseparator

% --- Repeated Games ---
\subsubsection{Repeated Games}

\begin{defbox}[Repeated Game]
A \textbf{repeated game} consists of the same \textbf{stage game} (in normal form) played multiple times by the same set of players.

Key questions:
\begin{itemize}
  \item Do agents observe what others played in previous rounds?
  \item Do they remember their own history?
  \item How is the overall utility defined from the sequence of stage-game payoffs?
\end{itemize}
\end{defbox}

\begin{defbox}[Strategies in Repeated Games]
\begin{itemize}
  \item A \textbf{stationary strategy}: play the same action in every round (memoryless)
  \item In general, the action at round $t$ can depend on the \textbf{entire history} of play up to $t$
  \item History-dependent strategies are crucial for sustaining cooperation in infinitely repeated games
\end{itemize}
\end{defbox}

\separator

% --- Finitely Repeated Games ---
\subsubsection{Finitely Repeated Games}

\begin{thmbox}[Backward Induction in Finitely Repeated Games]
In a finitely repeated game with a \textbf{known} number of rounds:
\begin{itemize}
  \item Can be represented as an extensive form game with imperfect information
  \item Payoffs are typically \textbf{additive} (sum of stage-game payoffs)
  \item At each stage, players observe what was played in previous rounds but not the current round
  \item Backward induction applies: in the last round, play the stage-game NE; knowing this, play the stage-game NE in the second-to-last round; and so on
\end{itemize}

\textbf{Result:} The only NE of a finitely repeated Prisoner's Dilemma is to \textbf{always Defect} in every round.
\end{thmbox}

\begin{tipbox}[Exam Tip: Finite Repetition Does Not Help]
A common exam question: ``Can cooperation be sustained in a finitely repeated Prisoner's Dilemma?'' The answer is \textbf{no}---backward induction unravels cooperation from the last round.
\end{tipbox}

\separator

% --- Infinitely Repeated Games ---
\subsubsection{Infinitely Repeated Games}

\begin{defbox}[Payoff Definitions for Infinite Games]
Since the game never ends, payoffs cannot be attached to terminal nodes. Two standard definitions:

\textbf{1.\ Average reward:}
\[
\bar{u}_i = \lim_{k \to \infty} \frac{\sum_{j=1}^{k} u_i^{(j)}}{k}
\]

\textbf{2.\ Future discounted reward} (more common):
\[
U_i = \sum_{j=1}^{\infty} \beta^{j} \cdot u_i^{(j)}
\]
where $\beta \in [0,1)$ is the \textbf{discount factor}.
\end{defbox}

\begin{defbox}[Discount Factor $\beta$: Two Interpretations]
\begin{enumerate}
  \item \textbf{Time preference:} The agent cares more about near-term payoffs than distant ones.
  \item \textbf{Continuation probability:} The game continues to the next round with probability $\beta$; it stops with probability $1 - \beta$.
\end{enumerate}

A higher $\beta$ (closer to 1) means the agent is more \textbf{patient} / the game is more likely to continue.
\end{defbox}

\separator

% --- Cooperation Strategies ---
\subsubsection{Sustaining Cooperation: Tit-for-Tat and Trigger}

\begin{defbox}[Tit-for-Tat (TFT)]
\begin{enumerate}
  \item Start by \textbf{cooperating}
  \item In round $j+1$, play whatever the opponent played in round $j$
\end{enumerate}

\textbf{Properties:} Forgiving (returns to cooperation after opponent cooperates again), simple, retaliatory.
\end{defbox}

\begin{defbox}[Trigger Strategy (Grim Trigger)]
\begin{enumerate}
  \item Start by \textbf{cooperating}
  \item If the opponent \textbf{ever} defects, defect \textbf{forever}
\end{enumerate}

\textbf{Properties:} Unforgiving (permanent punishment), strong deterrent, harsher than TFT.
\end{defbox}

\begin{thmbox}[Cooperation in Infinitely Repeated PD]
Both Tit-for-Tat and Trigger strategy are \textbf{Nash equilibria} in the infinitely repeated Prisoner's Dilemma for \textbf{sufficiently large} discount factor $\beta$.

Moreover, they form an equilibrium not just with themselves but also \textbf{with each other}---if one player uses TFT and the other uses Trigger, mutual cooperation is sustained.

\textbf{Intuition:} The short-term gain from defecting is outweighed by the long-run punishment, provided the player cares enough about future payoffs ($\beta$ large enough).
\end{thmbox}

\begin{exbox}[Computing Minimum $\beta$ for Trigger Strategy]
Consider the Prisoner's Dilemma with payoffs:
\begin{center}
\renewcommand{\arraystretch}{1.2}
\begin{tabular}{@{} c | c c @{}}
 & C & D \\
\midrule
C & $-1, -1$ & $-4, 0$ \\
D & $0, -4$ & $-3, -3$ \\
\end{tabular}
\end{center}

\textbf{If both play Trigger (cooperate forever):}
\[
U_{\text{coop}} = \sum_{j=1}^{\infty} \beta^j (-1) = \frac{-\beta}{1 - \beta}
\]

\textbf{If player deviates (defect once, then both defect forever):}
\[
U_{\text{dev}} = \beta \cdot 0 + \sum_{j=2}^{\infty} \beta^j (-3) = \frac{-3\beta^2}{1 - \beta}
\]

\textbf{Cooperation is a NE when:} $U_{\text{coop}} \geq U_{\text{dev}}$
\[
\frac{-\beta}{1 - \beta} \geq \frac{-3\beta^2}{1 - \beta} \implies -1 \geq -3\beta \implies \beta \geq \frac{1}{3}
\]

So Trigger sustains cooperation for $\beta \geq \frac{1}{3}$.
\end{exbox}

\bigseparator

% --- Summary Table ---
\subsubsection{Summary: Key Concepts}

\begin{center}
\renewcommand{\arraystretch}{1.4}
\begin{tabular}{@{} l p{9.5cm} @{}}
\toprule
\textbf{Concept} & \textbf{Key Fact} \\
\midrule
Perfect-info extensive form & 8-tuple $(N, A, H, Z, \chi, \rho, \sigma, u)$; game tree with full history visible \\
Pure strategy (extensive) & Complete plan: one action per choice node (or info set), even unreachable ones \\
Induced normal form & Exponentially larger; many cells map to the same leaf \\
Perfect info $\Rightarrow$ pure NE & Always exists (backward induction) \\
SPE & NE in every subgame; eliminates non-credible threats \\
Backward induction & Depth-first, linear-time algorithm to compute SPE \\
Imperfect info & Information sets partition choice nodes; player cannot distinguish nodes in same set \\
Finitely repeated PD & Only NE is always Defect (backward induction unravels) \\
Infinitely repeated PD & Cooperation via TFT/Trigger for large enough $\beta$ \\
Discount factor $\beta$ & Time preference or continuation probability; higher $\beta$ = more patient \\
\bottomrule
\end{tabular}
\end{center}

\bigseparator
